% Scope (SPEC.md, Table of contents): supergraph and subgraphs, entity ownership, keys, the Apollo Federation v2 directive tour, and why this book builds on that specification rather than on the Composite Schemas draft
\chapter{The Federation Model}
\label{ch:federation-model}

Federation is a join performed by something that owns neither side, and it
works only because every type that crosses a seam publishes a key both sides
agree on.

Your service already performs that join. Chapter~\ref{ch:data-without-n-plus-1}
made it cheap and chapter~\ref{ch:schema-design} gave both ends of it a name,
and in neither chapter did the key ever leave the process. This chapter draws
the line the key has to cross, and everything the model asks of you follows
from that one crossing.

\section{The Join You Cannot See}
\label{sec:ch05-the-join-you-cannot-see}

\index{join!inside one process|(}%
Start the service and ask for the schedule with its speakers.

\begin{minted}{text}
POST /graphql HTTP/1.1
Host: localhost:5001
Content-Type: application/json

{"query":"{ sessions(first: 10) { nodes { title speaker { name } } } }"}
\end{minted}

\begin{minted}[fontsize=\footnotesize]{json}
{"data":{"sessions":{"nodes":[
  {"title":"Schemas That Outlive Their Authors","speaker":{"name":"Ada Fischer"}},
  {"title":"Reading a Query Plan","speaker":{"name":"Ada Fischer"}},
  {"title":"Paging Without Offsets","speaker":{"name":"Bruno Kaminski"}},
  {"title":"The Bank That Split Its Graph","speaker":{"name":"Chidi Okafor"}}]}}}
\end{minted}

\index{statement count!the seam query}%
Two statements, as chapter~\ref{ch:data-without-n-plus-1} left it. The second
one is the one to look at.

\begin{minted}{text}
-- statement 2
SELECT "s"."Id", "s"."Bio", "s"."Name"
FROM "Speakers" AS "s"
WHERE "s"."Id" IN (@ids1, @ids2, @ids3)
\end{minted}

\index{key!carried by the foreign key}%
Three parameters for four sessions, because Ada Fischer gives two of them.
Those three values are the join. They came off the \code{SpeakerId} column of
the four rows the first statement returned, went into a DataLoader, and came
back out as a set. Every one of them is an integer that lived for the length of
one request inside one process.

\index{response!carries no key}%
Now read the response again and look for them. They are not there. Neither is
anything else that ties a speaker to a session: no \code{speakerId}, no
\code{sessionId}, no key of any kind. The nesting is the whole record of the
join, and a client that flattened this document into two lists would destroy
information it could not recover.

\index{SpeakerId@\code{SpeakerId}!absent from the schema by design}%
That is not an accident of what this query asked for. There is no field
anywhere on \code{Session} that would return it.
Chapter~\ref{ch:data-without-n-plus-1} put \code{[GraphQLIgnore]} on
\code{Session.SpeakerId} on purpose, and chapter~\ref{ch:schema-design} took
the other integer key out of the schema as well. A reader who typed this
service out of the book has, deliberately, no way to ask it which speaker
belongs to which session except by asking for the speaker.

\index{join!is a private implementation detail}%
So the join is real, it is cheap, and it is invisible from outside. It is also,
right now, a private matter between \code{SessionType.cs} and a table. Nothing
in the schema commits you to performing it, nothing in the schema describes
how, and you could swap the DataLoader for a dictionary this afternoon without
telling anyone.

\index{seam!where the line goes}%
Draw the line. Speakers belong to another team now: another repository, another
deployment, another database that your process cannot open a connection to.
The first statement is still yours. The second one is not, and the resolver
that would have issued it is holding three integers with nothing to do.

\index{join!has to move outward}%
There are only three places that join can go. It can go into every client,
which is the answer chapter~\ref{ch:do-you-need-this} spent a chapter arguing
against, and which the schema you have just built cannot even support. It can
stay in your service, which means your service calls theirs, which is a single
GraphQL service with extra steps and the same coordination problem underneath
it. Or it can move outward, into something in front of both services that
belongs to neither.

Figure~\ref{fig:ch05-the-join-moves} is the third answer, drawn beside the
arrangement you have now.

\begin{figure}
  \centering
  % The same join, performed twice: once inside one process, once across a seam.
%
% Same idiom as figures/tikz/ch01-one-field.tex, ch02-two-requests.tex and
% ch03-one-batch.tex: each lane is a timeline read left to right, ticks mark
% stages, a dashed vertical marks a boundary the work crosses, and a brace
% under the stages that matter carries a one-line label. Lane A is the
% arrangement chapter 3 finished with; lane B is the same join with a router
% in the middle. The four ticks sit at the same x in both lanes, because the
% argument of the figure is that the two arrangements do the same work in the
% same order and differ only in what the keys have to travel through, so any
% horizontal offset between the lanes would be saying something false.
%
% This figure fixes two visuals SPEC decision 33 says are reused. The seam is
% a dashed vertical rule; one round trip out and back crosses it twice, so
% lane B carries two rules and not four. The router is drawn as the actor
% whose lane it is, named in the tick labels, rather than as a box, because a
% box would make it look like a fourth service rather than the thing holding
% the plan.
%
% Bare tikzpicture (house style): the figure environment, caption and label
% live at the call site in the section file.
\begin{tikzpicture}[
  x=1mm, y=1mm,
  every node/.append style={font=\sffamily\scriptsize},
  axis/.style={draw, line width=0.4pt, -{Straight Barb[length=1.4mm, width=1.6mm]}},
  tick/.style={draw, line width=0.4pt},
  event/.style={anchor=south, align=center, text width=24mm, inner sep=1pt},
  span/.style={draw, line width=0.4pt},
  spanlabel/.style={anchor=north, align=center, inner sep=2pt},
  lane/.style={anchor=west, font=\sffamily\scriptsize\itshape},
  boundary/.style={draw, line width=0.4pt, dashed},
]

% ------------------------------------------------------------------- lane A
\node[lane] at (0,14) {A. one service, one process: the join you have now};

\draw[axis] (0,0) -- (150,0);
\foreach \x in {16,58,100,138}{\draw[tick] (\x,-1.5) -- (\x,1.5);}

\node[event, text width=26mm] at (16,3)  {sessions resolver:\\statement 1, 4 sessions};
\node[event, text width=26mm] at (58,3)  {speaker resolver x4:\\3 distinct keys held};
\node[event, text width=26mm] at (100,3) {statement 2: the 3\\keys in one IN list};
\node[event, text width=20mm] at (138,3) {one response\\assembled};

\draw[span] (50,-4) -- (50,-7) -- (108,-7) -- (108,-4);
\node[spanlabel] at (79,-8) {three integers, held in memory,\\named in no schema anywhere};

% ------------------------------------------------------------------- lane B
\begin{scope}[shift={(0,-50)}]
  \node[lane] at (0,14) {B. two subgraphs and a router: the same join, the same order};

  \draw[axis] (0,0) -- (150,0);
  \foreach \x in {16,58,100,138}{\draw[tick] (\x,-1.5) -- (\x,1.5);}
  \foreach \x in {37,119}{
    \draw[boundary] (\x,-2) -- (\x,8);
    \node[spanlabel, inner sep=1pt] at (\x,-2) {seam};
  }

  \node[event, text width=26mm] at (16,3)  {router asks sessions\\for the page};
  \node[event, text width=26mm] at (58,3)  {sessions answers: a\\speaker key per session};
  \node[event, text width=26mm] at (100,3) {router asks speakers,\\the 3 keys in one call};
  \node[event, text width=20mm] at (138,3) {router merges,\\one response};

  \draw[span] (50,-6) -- (50,-9) -- (108,-9) -- (108,-6);
  \node[spanlabel] at (79,-10) {the same three keys, serialized, through\\a field both services have to declare};
\end{scope}

\end{tikzpicture}

  \caption{The same join, twice. Lane A is the service as
  chapter~\ref{ch:data-without-n-plus-1} left it: the three keys are collected
  and spent without ever being named in a schema. Lane B does the same work in
  the same order, with two round trips through the middle of it. Nothing about
  the join changed, and that is the argument: what changed is that the keys now
  have to be serialized, which means they have to be fields, which means both
  services have to agree on what they are called and what is in them.}
  \label{fig:ch05-the-join-moves}
\end{figure}

\index{key!has to become public}%
Read the two lanes for what is different, which is not the number of steps. In
lane A three integers move between two lines of C\# in the same call stack. In
lane B the same three values have to survive being written to JSON, sent over a
network, parsed by a service that has never seen your database, and matched
against rows in a table you do not own. A value that has to do all of that is
not an implementation detail any more. It is an API.

\index{key!is the whole of the model}%
So the rest of this chapter is about who is allowed to be on each side of that
boundary, and what each of them has to declare before the value may cross.
\index{join!inside one process|)}%

\section{One Endpoint, Many Owners}
\label{sec:ch05-one-endpoint-many-owners}

\index{subgraph!is an ordinary service|(}%
A subgraph is your service with four additions. It still serves its own schema
at its own address, it still answers ordinary queries from anyone who asks, and
you can still open it in an IDE and browse it as chapter~\ref{ch:hot-chocolate-zero}
did. Chapter~\ref{ch:do-you-need-this} counted the four as a cost, and they are;
here is what each one buys.

\index{link@\code{"@link}!opts a schema into federation}%
\index{service@\code{\_service}!publishes the schema}%
\index{entities@\code{\_entities}!the entity route}%
The schema applies \code{@link} to itself, which is the opt-in: it says which
version of the federation vocabulary this schema is written in and which
directives it uses. It defines those directives, so that a composer reading the
schema alone knows what they mean. It resolves \code{\_service}, a field whose
one job is to hand back the schema as a string, so the composer can ask the
running service what it now looks like instead of being sent a file. And it
resolves \code{\_entities}, which is the entity route: a root field that takes
keys and returns the objects those keys name.

\index{subgraph!the additions are public}%
None of those four is hidden. They are fields on \code{Query} and directives in
the schema, visible to anybody who introspects the service. That is the fact
chapter~\ref{ch:entities-authorization} is built on.

\subsection{The Schema Nobody Writes}

\index{supergraph!is a composed result}%
Take three subgraph schemas and run them through a composer, and what comes out
is one schema in which \code{Session} carries the fields Sessions declared and
the fields Ratings declared, as though a single service had always had them.
That result is the supergraph. It is not a file anyone maintains, it is not
checked into a repository as a source document, and no team owns it, because it
is derived from the three that are owned.

\index{supergraph!is not a specification term}%
It is also not a word the specification defines. The Apollo Federation subgraph
specification regulates exactly one thing, the subgraph: what it must expose,
what those fields must return, and what the directives mean~\autocite{apollosubgraphspec}.
\enquote{Supergraph} appears in Apollo's own conceptual documentation, under a
heading that introduces it as terminology to learn, defined as \enquote{the
single, federated graph} you get \enquote{when combining multiple GraphQL
APIs}~\autocite{apollofederatedschemas}. That is a useful word and I use it
throughout this book. Which of the two has normative text behind it matters
when two implementations disagree, because the subgraph contract is the only
one of them you can hold anybody to.

\index{router!serves the supergraph}%
The router serves the supergraph, and it is the only address a client has. A
client sends it an ordinary GraphQL operation against an ordinary-looking
schema and gets an ordinary response. Nothing in that response says which
service produced which field, and no part of the protocol offers to tell the
client. That is the design working. It is also
chapter~\ref{ch:blame-routing}'s entire problem, arriving early.

\subsection{Two Different Times, Two Different Failures}

\index{composition!happens before any request}%
The join in figure~\ref{fig:ch05-the-join-moves} needs somebody to have worked
out, in advance, that a session's speaker is answered by asking the Speakers
service for an entity. Working that out is composition, and it happens before a
single request has been served. It reads the three schemas, checks that every field is
reachable from the root through some sequence of services, and produces
something the router can execute against.

\index{composition!fails on somebody else's change}%
Composition fails on schemas, not on data. It fails after everyone has merged,
it fails for the graph rather than for one service, and it can fail because of
a change you did not make. Chapter~\ref{ch:composition} is where you run it,
chapter~\ref{ch:satisfiability} is about the errors that are genuinely hard to
read, and appendix~\ref{app:comperrors} decodes the rest.

\index{execution!happens per request}%
Execution is the other time. A query arrives, the router works out which
services to ask in which order, asks them, and merges what comes back. Its
failures are the ones any distributed system has: a service that is slow, a
service that is down, a field that resolves to null and takes a branch of the
response with it. Part~III is almost entirely
about those.

\index{failure!which time it belongs to}%
Sorting a failure into the right one of those two is, in my experience, most of
debugging a federated graph, and the sort is cheap once you think to make it. A
field that has not worked since the deployment, for anyone, is composition. A
field that works for some requests and not others, or worked an hour ago, is
execution. The two have separate tooling, separate logs and separate people to
ask, so the cost of guessing wrong is an afternoon.
\index{subgraph!is an ordinary service|)}%

\input{chapters/05-the-federation-model/03-an-entity-is-a-type-with-a-key}
\section{Who Owns a Field}
\label{sec:ch05-who-owns-a-field}

\index{ownership!the rule|(}%
The rule is one sentence. A field belongs to the subgraph that owns the data
behind it, and no other subgraph declares it.

\index{ownership!an edge and its target are owned separately}%
Apply it to \code{Session.speaker} and it splits in a way that is easy to get
backwards. The field is on \code{Session}, and which speaker gives a session is
a fact the Sessions service stores, in a column, next to the row. So Sessions
owns \code{speaker}. What Sessions does not own is the speaker: the name, the
biography, everything the field's type promises. Sessions declares the field,
produces the key, and stops there. The edge and its target have different
owners, and the seam runs between them rather than through either.

\index{ownership!a subgraph contributing to a type it does not own}%
The second arrangement is the one that makes federation more than remote
procedure calls. The Ratings service stores nobody's sessions and nobody's
speakers, and it is nevertheless the right place for a session's ratings to
live. So it declares \code{Session} too, carrying the key and two of the three
fields chapter~\ref{ch:second-third-subgraph} gives it; the third needs a
directive section~\ref{sec:ch05-the-directives} has not introduced yet:

\begin{graphqlsdl}
type Session @key(fields: "id") {
  id: ID!
  averageScore: Float
  ratingCount: Int!
}
\end{graphqlsdl}

\index{Session@\code{Session}!assembled from two services}%
Compose that with the Sessions subgraph and a client sees one \code{Session}
carrying eight fields, two of which come from a service that stores no
sessions, and nothing in the schema says which came from where. A client asking
for a title and an average score has asked one question of one endpoint. That
is what the architecture is for, and it is delivered by a type appearing twice
with different fields on it.

\index{ownership!is a claim the composer enforces}%
Which is why ownership is not a convention here. It is a claim each subgraph
makes in its schema, and the composer checks the claims against each other
before anything runs. Chapter~\ref{ch:do-you-need-this} argued that federation
does not remove coordination between teams but moves it into what a program
checks. This is that program, and these are the checks.

\subsection{When Two Subgraphs Claim the Same Field}

\index{shareable@\code{"@shareable}!declares a field may be duplicated}%
Declare a field in two subgraphs and composition fails. Not with a warning, not
with a last-one-wins: it fails, and the message names the field and every
subgraph that declared it, and says that none of them declared it
\code{@shareable}. The composer's position is that a field defined twice is
either a mistake or a decision, and it declines to guess which.

\index{shareable@\code{"@shareable}!is a promise about the value}%
\code{@shareable} is how you say it was a decision. It means the field may be
resolved by more than one subgraph and that each of them will return the same
value for the same object, which is a promise about your data and not a
compiler flag. Where that promise holds, the router gains a choice: it can get
the field from whichever subgraph it was already talking to and save a hop.
Where it does not hold, you have built a graph that returns different answers
depending on a routing decision no client can see.

\index{node@\code{node}!collides across subgraphs}%
This book has already bought itself one of these, and the place to meet it is
here rather than in the middle of chapter~\ref{ch:first-subgraph}.
Chapter~\ref{ch:schema-design} turned on global object identification, which
puts \code{node} and \code{nodes} on \code{Query}. Every subgraph in this book
will therefore export both. Two subgraphs, two \code{Query.node} fields, and
composition stops with exactly the error above. The fix is to declare them
shareable in both, and the interesting part is that the promise is real: any
subgraph handed a global node id has to answer with the same object, which is
the promise chapter~\ref{ch:schema-design}'s id format already makes.

\subsection{Moving a Field Between Owners}

\index{override@\code{"@override}!moves a field}%
Ownership changes. A field that started in the service that had the data ends
up in the service that grew up around it, and the change has to happen without
a coordinated deployment, because avoiding coordinated deployments is why you
federated. \code{@override(from:)} is the mechanism: the receiving subgraph
declares the field and names the subgraph it is taking it from, composition
gives it the field, and the old subgraph keeps its copy until somebody deletes
it.

\index{override@\code{"@override}!two deployments, either order}%
That ordering is the point. Both subgraphs carry the field for as long as you
like, so the two deployments happen on two teams' calendars in either order,
and the graph is answerable throughout. It is the same problem
chapter~\ref{ch:schema-design} solved for a field a client depends on, one level
up, with services in the place of clients.

\index{ownership!the part no directive settles}%
None of which decides who should own a field, and that is not an oversight in
the specification. \code{@shareable} lets two teams both be right;
\code{@override} lets one take a field from another; neither has an opinion
about whether they should. Chapter~\ref{ch:ownership} is the chapter about the
part that stays a conversation, and I put it in part~III with the failures
rather than here with the mechanisms because that is where I think it belongs.
\index{ownership!the rule|)}%

\section{The Directives, and What Each One Is For}
\label{sec:ch05-the-directives}

\index{directives!the vocabulary|(}%
Federation v2 hands every subgraph library a block of sixteen directive
definitions to add to your schema for you, plus one more for libraries that
cannot write GraphQL's own \code{extend}
keyword~\autocite{apollosubgraphspec}. Four of the sixteen carry the model.
The rest get a paragraph each so that you recognize one in somebody else's
schema, and appendix~\ref{app:directives} carries the definitions in full.

\subsection{The Opt-In}

\index{link@\code{"@link}!is the version declaration}%
\code{@link} comes first because nothing else exists without it. It goes on the
schema itself and says which version of the vocabulary this document is written
in and which of its names the document uses:

\begin{graphqlsdl}
directive @link(url: String!, as: String, for: link__Purpose,
                import: [link__Import]) repeatable on SCHEMA
\end{graphqlsdl}

\index{link@\code{"@link}!the import list}%
The url ends in a version, and a subgraph that names an older one gets the
vocabulary of that older one, which is how a graph composes while its services
are on different versions. The import list does the other job: a directive you
did not import is not available under its bare name. In practice the library
writes this line for you from what your schema actually uses, and the first
time you will care is the day you reach for a directive your version does not
have and get an error about an unknown name rather than about a version.

\index{link@\code{"@link}!Hot Chocolate ships a narrower one}%
One caution. The definition above is the specification's. The one Hot
Chocolate's federation package writes into an exported schema has neither
\code{as} nor \code{for}, and types \code{import} as a plain list of strings, so
the namespacing \code{as} exists for is not available to you here. Nothing in
this book needs it, and I say so because reading the specification and
then reaching for that argument is a natural thing to do.

\subsection{The Seam}

\index{directives!the four that describe a seam}%
Four directives describe what crosses a boundary, and
section~\ref{sec:ch05-an-entity-is-a-type-with-a-key} has already spent itself
on the first. \code{@key} makes a type an entity. The other three qualify what
a subgraph may say about an entity it does not own:

\begin{graphqlsdl}
directive @external on FIELD_DEFINITION | OBJECT
directive @requires(fields: FieldSet!) on FIELD_DEFINITION
directive @provides(fields: FieldSet!) on FIELD_DEFINITION
\end{graphqlsdl}

\index{external@\code{"@external}!declares a field this subgraph does not own}%
\code{@external} marks a field that appears in this subgraph's schema but is
owned somewhere else. It is a declaration of borrowing: the field is here so
that something in this schema can refer to it, and this service does not
resolve it.

\index{requires@\code{"@requires}!asks the router for data first}%
\code{@requires} is what borrowing is usually for. A field you own can need a
value you do not have, and \code{@requires} names that value as a field set, so
the router fetches it from its owner and hands it to your resolver before your
resolver runs. It turns one fetch into two that must happen in order, and that
ordering is the price of the convenience.

\index{provides@\code{"@provides}!offers a field the router would otherwise fetch}%
\code{@provides} runs the other way. It says that along this particular field,
this subgraph can already supply some of the target entity's fields, so the
router may skip a hop it was going to make. It is an optimization with a
promise attached, and the promise is the same one \code{@shareable} makes: the
value has to match what the owner would have said.

Chapter~\ref{ch:second-third-subgraph} is where the first two get used on a
real seam and the third is weighed and left out, and
chapter~\ref{ch:satisfiability} is where getting them wrong produces
the composition errors that are hardest to read.

\subsection{Overlap and Visibility}

\index{directives!the four that manage overlap}%
\code{@shareable} and \code{@override} are
section~\ref{sec:ch05-who-owns-a-field}'s, and two more belong beside them
because they also change what the supergraph shows rather than what a subgraph
computes.

\index{inaccessible@\code{"@inaccessible}!hides a field from clients}%
\code{@inaccessible} takes a field out of the client's reach without taking it
out of the graph. Apollo is precise about the distinction and it is one to
hold on to: the field \enquote{is not omitted from the supergraph schema, so
the router still knows it exists (but clients can't include it in
operations)}~\autocite{apollodirectives}. Other subgraphs may still reference
it. It is the way to add a field to one service before the others are ready to
have it exist, and the way to take one off the public graph without taking it
out of the code.

\index{tag@\code{"@tag}!labels a field for tooling}%
\code{@tag} \enquote{applies arbitrary string metadata to a schema
location}~\autocite{apollodirectives} and means nothing by itself. What reads
the tags is whatever you point at the graph: a tool that builds a restricted
view of the supergraph for one audience, a policy check, a report. It is a
hook, and the value comes from what you hang on it.

\subsection{Access Control}

\index{directives!the three that carry authorization}%
\code{@authenticated}, \code{@requiresScopes} and \code{@policy} let a subgraph
state a requirement in its schema as well as in its resolvers, so that the
router can enforce it too. Chapter~\ref{ch:entities-authorization}
is about the reason to care, which is not the one people expect: the entity
route is a second door into every object in your graph, and a guard written on
the field a client normally uses does not stand in front of it.

\subsection{The Rest}

\index{interfaceObject@\code{"@interfaceObject}!contributes to every implementation}%
\code{@interfaceObject} lets a subgraph add a field to every type implementing
an interface without knowing what those types are. Apollo describes the
subgraph's declaration as \enquote{an abstraction of another subgraph's entity
interface}~\autocite{apollodirectives}: you declare an object type carrying the
interface's name and its key, and composition distributes the field to each
implementation.

\index{composeDirective@\code{"@composeDirective}!carries a custom directive through}%
\code{@composeDirective} names one of your own directives and asks the composer
to carry it into the supergraph. It is needed because \enquote{by default,
composition omits most directives from the supergraph
schema}~\autocite{apollodirectives}, which is the right default and an awkward
surprise the first time a directive of yours vanishes.

\index{extends@\code{"@extends}!for libraries without the keyword}%
\code{@extends} is the seventeenth, and it exists only for libraries that
cannot write GraphQL's own \code{extend} keyword. Hot Chocolate can, so this
book never uses it, and it is here because you will meet it in schemas written
by libraries that cannot.

\index{directives!what this package does not ship}%
Two of the sixteen are not in the package this book uses. \code{@context} and
\code{@fromContext}, which pass a value down from an ancestor field into a
resolver's argument, have no attribute and no directive type in
\code{HotChocolate.ApolloFederation}, so no chapter here claims anything about
them.

\index{cost@\code{"@cost}!arrives from elsewhere}%
\code{@cost} and \code{@listSize} are a different case again. Apollo documents
both alongside the sixteen, as controls on how much work a client may ask
for rather than as anything to do with a seam, and Hot Chocolate implements
them in its core rather than in the federation package. Which is why
chapter~\ref{ch:data-without-n-plus-1}'s exported schema was already carrying
both before this book had said the word subgraph.
\index{directives!the vocabulary|)}%

\section{Which Federation, and Why}
\label{sec:ch05-which-federation-and-why}

\index{Composite Schemas specification|(}%
Chapter~\ref{ch:do-you-need-this} said that federation had outgrown the company
that named it, and that the GraphQL Specification Working Group had a
subcommittee producing a common specification for it. It then said this book
builds on Apollo Federation v2 anyway. Here is the reasoning, because a reader
who has just been walked through one directive vocabulary and is about to
spend the rest of a book on it is owed it.

\index{Composite Schemas specification!has no finished edition}%
The first reason is that there is nothing finished to build on. The GraphQL
Composite Schemas Spec publishes exactly one edition, labeled \enquote{Working
Draft}, and the project as a whole is marked \enquote{Prerelease}. There is no
numbered release and never has been~\autocite{compositeschemasspec}. That is not
a criticism of the work. It is a statement about what a book can responsibly
teach: a draft that may change is a fine thing to follow and a poor thing to
print several hundred pages against.

\index{Composite Schemas specification!is not a renaming}%
The second reason is that it is not Apollo Federation v2 with the serial
numbers filed off, so the choice is real rather than cosmetic. Two differences
are enough to make the point. Apollo's key is a \code{FieldSet}; the draft's is
a \code{FieldSelectionSet}, a different scalar with its own grammar. And where
Apollo puts a reserved root field on every subgraph, the draft marks an ordinary
field:

\begin{graphqlsdl}
directive @lookup on FIELD_DEFINITION
\end{graphqlsdl}

\index{entities@\code{\_entities}!absent from the draft}%
There is no \code{\_entities} in that document at all. A subgraph written to the
draft says \enquote{this existing field can be used to look an entity up}, where
a subgraph written to Apollo Federation says \enquote{here is a reserved field
that takes representations}. Those are different designs with different failure
modes, and half of part~III of this book is about the failure modes of the
second one.

\index{Fusion!cannot span both versions}%
The third reason is the one that actually decided it. The .NET implementation
of the draft is Fusion, and Fusion 16 is a rewrite rather than an upgrade.
Michael Staib, who founded ChilliCream, puts it plainly on their blog:
\enquote{When we set out to rewrite Fusion, we had three main goals: remove the
constraints imposed by Hot Chocolate, achieve top-tier performance with .NET,
and make the gateway feel like a natural extension of your ASP.NET Core app,
rather than a black box configured with YAML}~\autocite{staib2026fusion16}. The
same post says Fusion 16 implements the current version of the Composite Schema
Specification.

\index{HotChocolate.ApolloFederation@\code{HotChocolate.ApolloFederation}!spans three majors}%
A reader still on Hot Chocolate~14 cannot follow any of that, because the
Fusion that shipped in the 14 era is a different product that stopped receiving
releases and was replaced by differently named packages. \code{HotChocolate.ApolloFederation}
was not. It has an unbroken release line across all three major versions this
book could have covered, which makes Apollo Federation v2 the only federation
contract a reader on either of the two it does cover can write to. That is the
decision, and it is one about readers rather than about the merits of two
specifications.

\index{Composite Schemas specification!extends rather than replaces}%
None of which is a bet against the draft. The working group's own summary of its
August 2026 meeting records two things: that the specification's
final name will be the GraphQL Federation Specification, and that it
\enquote{would extend Apollo Federation rather than replace
it}~\autocite{compositeschemaswg202608}. I attach a caveat to that source rather
than to the claim, because the document carries one itself: it is an
automatically generated meeting summary and says so at the top. The rename had
not reached the specification when I checked on 23 August 2026: the published
document still called itself the Composite Schemas Spec.

\index{Apollo Federation!is the thing to learn either way}%
So the vocabulary in this chapter is worth learning on either outcome. If the
draft lands as an extension of what Apollo published, the model you are about to
read is the base it extends. If it diverges further, you will have run a
federated graph in production and will be in a position to judge the difference,
which is more than you can do from either specification's text.
\index{Composite Schemas specification|)}%

\section{The Graph This Book Builds}
\label{sec:ch05-the-graph-this-book-builds}

\index{conference example!split into three|(}%
Three services, and the split is the obvious one: whoever stores the rows owns
the type. Sessions keeps what chapter~\ref{ch:schema-design} left it. Speakers
takes the other table. Ratings is new, and it exists because a service that
contributes to a type it does not own is the arrangement this model is for, and
a two-service example never shows it.

\index{Sessions subgraph!what it declares}%
Here is the Sessions subgraph, elided as in
chapter~\ref{ch:schema-design} down to what the seam touches. Gone from it:
the \code{Query} and \code{Mutation} types and their fields, the paging types,
the \code{Node} interface, and the descriptions and \code{@cost} directives the
export writes on everything.

\begin{graphqlsdl}
schema
  @link(
    url: "https://specs.apollo.dev/federation/v2.6"
    import: ["@key", "@tag", "FieldSet"]
  ) {
  query: Query
}

type Session @key(fields: "id") {
  id: ID!
  title: String!
  abstract: String
  startsAt: DateTime!
  durationMinutes: Int!
  speaker: Speaker
}

type Speaker @key(fields: "id") {
  id: ID!
}
\end{graphqlsdl}

\index{link@\code{"@link}!the library writes it}%
The header is the library's work rather than yours, and you read it once and
then stop seeing it. It names the version of the vocabulary this document is written in, and
it imports what the document uses. The last type in the document is the whole
cost of the seam on this side: one field, so that the key has somewhere to
live.

\index{Speakers subgraph!what it declares}%
The Speakers subgraph opens with the same header and declares the same type,
and means it:

\begin{graphqlsdl}
type Speaker @key(fields: "id") {
  id: ID!
  name: String!
  bio: String
}
\end{graphqlsdl}

\index{Ratings subgraph!contributes to a type it does not own}%
And Ratings declares a type it owns and a type it does not:

\begin{graphqlsdl}
type Rating @key(fields: "id") {
  id: ID!
  score: Int!
  comment: String
  session: Session!
}

type Session @key(fields: "id") {
  id: ID!
  ratings: [Rating!]!
  averageScore: Float
}
\end{graphqlsdl}

\index{Session@\code{Session}!declared twice}%
\code{Session} now appears in two of the three documents and \code{Speaker} in
two, and in each pair one declaration is the type and the other is the key. Read
down the three and you can see the whole graph: two seams, both of them at a
\code{@key(fields: "id")}, and Sessions sitting on both.

\subsection{What the Client Gets}

\index{supergraph!hides which service answered}%
Compose the three and \code{Session} arrives as one type:

\begin{graphqlsdl}
type Session {
  id: ID!
  title: String!
  abstract: String
  startsAt: DateTime!
  durationMinutes: Int!
  speaker: Speaker
  ratings: [Rating!]!
  averageScore: Float
}
\end{graphqlsdl}

\index{supergraph!carries no federation directives}%
No \code{@key}. No second declaration. Nothing saying that two of those fields
are answered by a service which has never heard of a conference schedule. The
directives did their work at composition and are gone from what the client
reads, which is the correct outcome and the reason
chapter~\ref{ch:blame-routing} exists: a client that cannot see the seam cannot
tell you which side of it went wrong.

\index{seam!is a design decision, not a discovery}%
Look at where the two seams are, because you chose them and you could have
chosen otherwise. Sessions could have kept a speaker's name and saved a hop.
Ratings could have been a field on the Sessions service and saved a service. In
both cases the reason not to is the one chapter~\ref{ch:do-you-need-this}
started with: a seam is where one team stops needing another team's permission,
and a seam drawn anywhere else is a network call bought for nothing.

\index{seam!the bill for this one}%
The bill for this particular decomposition arrives in
chapter~\ref{ch:cross-seam-paging}. Sorting the schedule by average score is
an obvious thing for a client to ask and there is no service that can answer
it: Sessions owns the ordering and knows no scores, Ratings owns the scores
and knows no schedule, and no directive in
section~\ref{sec:ch05-the-directives} closes that gap. I chose this example
partly because that problem is unavoidable in it rather than contrived.

\subsection{What Part II Does With It}

\index{Part II!the four chapters that build this}%
Chapter~\ref{ch:first-subgraph} makes the service you already have into a
subgraph: the package, the key, the reference resolver, and the two fields the
contract adds to \code{Query}. It gets most of the way to the first of those
three documents and stops short of one thing. Sessions still stores the
speakers, so its \code{Speaker} is the whole type rather than the stub above,
and it stays that way until there is a service to move the rows to. That
chapter ships one subgraph and composes nothing.

Chapter~\ref{ch:composition} composes, which means meeting the composer's
opinions about the schema you wrote, including the one
section~\ref{sec:ch05-who-owns-a-field} promised you about \code{Query.node}.
Chapter~\ref{ch:enter-the-router} puts a router in front and answers the first
query that crosses a seam. Chapter~\ref{ch:second-third-subgraph} adds the other
two services and the directives in
section~\ref{sec:ch05-the-directives}'s seam group, and ends with the graph
part~III spends nine chapters breaking.
\index{conference example!split into three|)}%

\begin{chaptersummary}
  \item Your service already performs the join federation performs. The
        difference is that three integers which lived inside one process now
        have to be serialized, which makes them fields, which makes them part
        of a contract two teams have to agree on. Everything in the model
        follows from that one crossing.
  \item A subgraph is an ordinary GraphQL service that also applies
        \code{@link} to its schema, defines the federation directives, and
        resolves \code{\_service} and \code{\_entities}. A supergraph is what a
        composer makes of several of them, and it is not a document anybody
        writes or owns. Only the first of those two words has a specification
        behind it.
  \item Composition and execution are different times with different failures.
        A field that has not worked since the deployment, for anybody, is
        composition. A field that works for some requests and not others is
        execution. They have separate tooling and separate people to ask.
  \item An entity is a type with a \code{@key}, and the key is a selection set
        written as a string. Declaring it forces every subgraph that touches
        the type to publish the key fields and to compute the same value for
        the same object, which nothing checks. \code{\_entities} takes those
        keys as untyped JSON and must answer in the order it was asked, because
        position is the only thing connecting an answer to its question.
  \item A field belongs to the subgraph that owns its data, and the composer
        enforces that as a claim rather than a convention: two subgraphs
        declaring one field fail composition unless both call it
        \code{@shareable}. Global object identification has already bought this
        book one of those, on \code{Query.node}.
\end{chaptersummary}

